\documentclass[11pt,a4paper]{article}

\usepackage[T1]{fontenc}
\usepackage[utf8]{inputenc}
\usepackage[french]{babel}
\usepackage{geometry}
\usepackage{enumitem}
\usepackage{hyperref}
\usepackage{titlesec}

\geometry{margin=2.3cm}
\setlength{\parskip}{0.4em}
\setlist[itemize]{noitemsep, topsep=2pt, leftmargin=1.2em}
\setlist[enumerate]{noitemsep, topsep=2pt, leftmargin=1.5em}

\titleformat{\section}{\normalfont\bfseries\large}{\thesection.}{0.5em}{}
\titleformat{\subsection}{\normalfont\bfseries\normalsize}{\thesubsection}{0.5em}{}

\hypersetup{colorlinks=true, linkcolor=black, urlcolor=blue}

\begin{document}

\begin{center}
    {\LARGE \textbf{Rapport d'exercice}}\\[0.6em]
    {\large \textbf{Gestion de Projet Informatique}}\\[0.4em]
    {\normalsize M1-GL / SIRI / SI --- IFRI}\\[0.3em]
    {\normalsize Année académique 2024--2025}
\end{center}

\vspace{0.5em}
\noindent\textbf{Objectif :} Rappels du cours

\noindent\textbf{Consigne :} Faire un résumé de deux à trois pages maximum pour résumer les notions essentielles vues dans le cours de Gestion de Projet Informatique jusqu'à cette date.

\noindent\textbf{Données :} Support de cours (PMBOK\textsuperscript{\textregistered} 6\textsuperscript{e} et 7\textsuperscript{e} éditions, PMI).

\vspace{0.8em}
\hrule
\vspace{0.8em}

\section{Projet, programme et portefeuille}

Un \textbf{projet} est une entreprise \textbf{temporaire} visant un \textbf{résultat unique} (produit ou service), avec un objectif à atteindre, des ressources limitées et un délai donné. Ses caractéristiques essentielles sont l'unicité du livrable, la progressivité, l'induction du changement et la génération d'un bénéfice quantifiable. À l'inverse, les activités opérationnelles récurrentes (production courante, maintenance) ne sont pas des projets.

Le \textbf{management de projet} est la démarche qui structure, assure et optimise le déroulement d'un projet. Il regroupe des techniques pour \textbf{planifier}, \textbf{estimer} et \textbf{contrôler} les activités afin de respecter délais, budget et périmètre. Chaque projet produit des \textbf{livrables} (résultats tangibles et mesurables) organisés en \textbf{phases} formant le \textbf{cycle de vie}.

Un \textbf{programme} regroupe des projets liés, gérés de façon coordonnée pour obtenir des bénéfices supérieurs à une gestion isolée. Un \textbf{portefeuille} rassemble projets et programmes au service des \textbf{objectifs stratégiques} de l'organisation.

Les phases peuvent être \textbf{séquentielles} (moins d'incertitude) ou se \textbf{chevaucher} (échéancier raccourci, risque accru). La \textbf{porte de phase} (GO/NO-GO) est une revue décisionnelle en fin de phase : poursuivre, modifier ou clore le projet. L'étude de faisabilité peut ainsi aboutir à une décision d'arrêt.

\section{Acteurs, documents et chef de projet}

Les principaux \textbf{acteurs} sont le \textbf{sponsor} (commanditaire, apporte ressources et légitimité), le \textbf{chef de projet}, l'\textbf{équipe projet}, le \textbf{comité de pilotage}, le \textbf{PMO}, les clients/utilisateurs, les acteurs techniques et les parties prenantes internes/externes.

Documents fondamentaux :
\begin{itemize}
    \item \textbf{Business case} : justification économique (TIR, VAN, ROI).
    \item \textbf{Charte de projet} : autorise le projet, définit objectifs, sponsor, équipe et niveau d'autorité du chef de projet.
    \item \textbf{Énoncé de périmètre} : base d'accord entre parties prenantes sur ce qui est inclus/exclu.
    \item \textbf{Exigences} : besoins quantifiés et documentés des parties prenantes.
\end{itemize}

Le \textbf{chef de projet} coordonne l'équipe comme un chef d'orchestre : il ne réalise pas toutes les tâches techniques mais assure vision, planification et aboutissement du projet. Il consacre environ \textbf{90\,\%} de son temps à la \textbf{communication}. Il équilibre contraintes divergentes (coût, délai, périmètre, qualité) et mobilise différents \textbf{pouvoirs} (position, expert, référence, récompense, coercitif, relationnel).

\textbf{Leadership $\neq$ management} : le leadership inspire et oriente ; le management organise et contrôle. Compétences clés : motivation, négociation, résolution de conflits, intelligence émotionnelle, facilitation. Outils d'équipe : objectifs \textbf{SMART}, \textbf{charte d'équipe}, \textbf{règles de base} (ground rules), théories de \textbf{Maslow} et \textbf{Herzberg}.

\section{Gestion des conflits}

Les conflits sont inévitables (rareté des ressources, priorités, styles de travail, différences culturelles). Leur gestion efficace améliore performance et productivité ; une gestion inefficace génère animosité et dégradation des résultats.

\textbf{Causes fréquentes} : concurrence pour les ressources, désaccords sur rôles et méthodes, différences d'objectifs/valeurs, problèmes de communication.

\textbf{Approches de résolution} (Thomas-Kilmann) :
\begin{itemize}
    \item \textbf{Évitement/retrait} : reporter ou se retirer.
    \item \textbf{Accommodement} : privilégier l'harmonie.
    \item \textbf{Compromis} : satisfaction partielle de chacun.
    \item \textbf{Forçage} : imposer une solution gagnant/perdant.
    \item \textbf{Collaboration} : recherche de consensus (approche privilégiée).
\end{itemize}

\section{PMI, PMBOK 6\textsuperscript{e} et 7\textsuperscript{e} édition}

Le \textbf{PMI} (Project Management Institute), fondé en 1969, publie les standards de référence et délivre des certifications dont le \textbf{PMP}. Le \textbf{PMBOK\textsuperscript{\textregistered}} recense les bonnes pratiques en gestion de projet.

\textbf{PMBOK 6} --- approche \textbf{processus} :
\begin{itemize}
    \item \textbf{10 domaines de connaissance} : intégration, périmètre, délais, coûts, qualité, ressources, communications, risques, approvisionnements, parties prenantes.
    \item \textbf{5 groupes de processus} : initialisation, planification, exécution, maîtrise, clôture.
\end{itemize}

\textbf{PMBOK 7} --- évolution vers les \textbf{principes} et la \textbf{valeur} :
\begin{itemize}
    \item \textbf{8 domaines de performance} : parties prenantes, équipe, approche de développement, planification, travaux, mesures, livraison, incertitude.
    \item \textbf{12 principes} : gestionnaire attentionné, équipe collaborative, implication des parties prenantes, focus sur la valeur, pensée holistique, leadership, sur-mesure (tailoring), qualité, complexité, risques, adaptabilité, conduite du changement.
    \item Passage des livrables aux \textbf{résultats à valeur ajoutée} ; prise en compte des approches \textbf{prédictives, agiles et hybrides}.
\end{itemize}

\section{Approches prédictive, agile et hybride}

\textbf{Prédictif (Waterfall)} : planification détaillée en amont, phases séquentielles, périmètre stable. Avantages : structure claire, estimation facilitée. Inconvénients : faible flexibilité, risque de produit obsolète si les besoins évoluent.

\textbf{Agile} : livraisons itératives, adaptation au changement, collaboration client. Le \textbf{Manifeste Agile} valorise les individus, le logiciel fonctionnel, la collaboration et la réactivité au changement.

\textbf{Scrum} (framework agile) :
\begin{itemize}
    \item \textbf{Rôles} : Product Owner (valeur et backlog), Scrum Master (facilitation), Développeurs (3 à 9).
    \item \textbf{Artéfacts} : Product Backlog, Sprint Backlog, Incrément (Definition of Done).
    \item \textbf{Événements} : Sprint (timebox $\leq$ 4 semaines), Daily Scrum (15 min), Sprint Planning, Sprint Review, Sprint Retrospective.
    \item \textbf{User story} : \textit{En tant que [QUI], je veux [QUOI] afin de [POURQUOI]}.
\end{itemize}

\textbf{Hybride} : combine planification prédictive et pratiques agiles (rétrospectives, engagement d'équipe) selon le contexte du projet.

\textbf{Cycles de vie} : prédictif, itératif, incrémental, agile, hybride --- choisis selon complexité, incertitude et stabilité des exigences.

\section{Organisation, PMO et environnement du projet}

\textbf{Structures organisationnelles} :
\begin{itemize}
    \item \textbf{Fonctionnelle} : autorité du chef de projet faible, loyauté au service.
    \item \textbf{Orientée projet} : autorité élevée du chef de projet, équipe dédiée.
    \item \textbf{Matricielle} : équilibre entre fonctionnel et projet (faible, équilibrée ou forte).
\end{itemize}

Le \textbf{PMO} (Project Management Office) centralise la gouvernance des projets. Types : \textbf{soutenant} (conseil, modèles), \textbf{de conformité} (contrôle des standards), \textbf{directif} (gestion directe des projets).

\textbf{EEF} (facteurs environnementaux) : éléments externes à l'équipe influençant le projet (culture, normes, infrastructure). \textbf{OPA} (actifs organisationnels) : ressources internes réutilisables (modèles, leçons apprises, données historiques, procédures).

\section{Code éthique et conduite professionnelle}

Le code du PMI s'applique aux membres et certifiés. Il repose sur quatre valeurs :
\begin{itemize}
    \item \textbf{Responsabilité} : assumer décisions, actions et conséquences.
    \item \textbf{Respect} : estime de soi, des autres et des ressources confiées.
    \item \textbf{Équité} : décisions impartiales, sans favoritisme.
    \item \textbf{Honnêteté} : vérité dans les communications et la conduite.
\end{itemize}

\vspace{0.6em}
\hrule
\vspace{0.4em}
\noindent\textit{Ce résumé synthétise les notions essentielles du module : fondamentaux du projet, rôle du chef de projet, gestion des conflits, référentiels PMI/PMBOK, approches prédictive et agile (Scrum), structures organisationnelles, PMO et code éthique professionnel.}

\end{document}
